\documentclass[12pt]{article}
\usepackage{amsmath,amssymb,amsthm}
\usepackage{fullpage}
\usepackage[T1]{fontenc}
\usepackage[utf8]{inputenc}
\usepackage{hyperref}

\newtheorem{theorem}{Theorem}
\newtheorem{lemma}[theorem]{Lemma}
\newtheorem{corollary}[theorem]{Corollary}
\newtheorem{proposition}[theorem]{Proposition}
\theoremstyle{definition}
\newtheorem{definition}[theorem]{Definition}
\theoremstyle{remark}
\newtheorem{remark}[theorem]{Remark}

\newcommand{\bR}{\mathbb{R}}
\newcommand{\om}{\omega}

\title{Solution to Problem 8:\\
Lagrangian Smoothings of Polyhedral Lagrangian Surfaces\\
with 4-Valent Vertices}
\author{}
\date{April 2026}

\begin{document}
\maketitle

\begin{abstract}
We prove that the answer is \textbf{yes}: every polyhedral Lagrangian surface
$K\subset\bR^4$ in which exactly 4 faces meet at every vertex necessarily
admits a Lagrangian smoothing.  The proof proceeds by analysing the local
model at each vertex (four Lagrangian half-planes meeting along a common
edge), showing that the local obstruction to smoothing vanishes under the
4-valent condition, and then applying a global gluing argument using the
Maslov index.
\end{abstract}

%------------------------------------------------------------------
\section{Setup and Definitions}

\subsection{Polyhedral Lagrangian surfaces}

Equip $\bR^4$ with the standard symplectic form $\om = dx_1\wedge dy_1 +
dx_2\wedge dy_2$.  A \emph{polyhedral Lagrangian surface} $K\subset\bR^4$
is a finite polyhedral 2-complex that is:
\begin{itemize}
\item A topological submanifold of $\bR^4$,
\item Each 2-face $\sigma$ is a Lagrangian polygon: $\om|_\sigma = 0$.
\end{itemize}

A \emph{Lagrangian smoothing} of $K$ is a Hamiltonian isotopy
$K_t\subset\bR^4$, $t\in(0,1]$, of smooth Lagrangian submanifolds,
extending to a topological isotopy $K_t$, $t\in[0,1]$, with $K_0 = K$.

\subsection{The 4-valent condition}

We assume: \emph{exactly 4 faces meet at every vertex of $K$}.
This is the analogue of a \emph{4-valent graph} in 2 dimensions.
Along each edge $e$, exactly 2 faces meet (as for a piecewise-linear
surface); the 4-valent condition restricts the combinatorial structure at
0-dimensional faces (vertices).

\subsection{Local model at a vertex}

Near a vertex $v$, the surface $K$ looks like a cone on a graph $\Gamma_v$
embedded in the link $S^3\subset\bR^4$.  The 4-valent condition says $\Gamma_v$
is a graph with all vertices of degree 4.  In the Lagrangian setting, the link
$\Gamma_v$ is a Legendrian graph in $(S^3, \xi_{\mathrm{std}})$ (the standard
contact structure on $S^3$).

The simplest 4-valent vertex model: four Lagrangian half-planes $P_1,P_2,P_3,P_4$
meeting at the origin, pairwise sharing an edge (a half-line through the origin).

%------------------------------------------------------------------
\section{Local Analysis at a Vertex}

\begin{lemma}[Local smoothing at a 4-valent vertex]\label{lem:local}
Let $v$ be a 4-valent vertex of a polyhedral Lagrangian surface $K$.  In a
neighbourhood $U$ of $v$ in $\bR^4$, there exists a Lagrangian smoothing of
$K\cap U$ relative to the boundary $\partial U$.
\end{lemma}

\begin{proof}
\textbf{Normal form.}
By a symplectomorphism of $(\bR^4,\om)$ (Darboux theorem applied locally),
we may assume $v = 0$ and that the four faces meeting at $v$ are locally
modelled as:
\begin{align*}
  F_1 &= \{(x_1,x_2,0,0) : x_1\geq 0,\, x_2\geq 0\}, \\
  F_2 &= \{(x_1,0,y_2,0) : x_1\leq 0,\, y_2\geq 0\}, \\
  F_3 &= \{(0,x_2,0,y_1) : x_2\leq 0,\, y_1\geq 0\}, \\
  F_4 &= \{(0,0,y_2,y_1) : y_2\leq 0,\, y_1\leq 0\}.
\end{align*}
(Each is a Lagrangian plane since $\om|_{F_i} = 0$: for $F_1$,
$\om|_{F_1} = dx_1\wedge dy_1 + dx_2\wedge dy_2 = 0$ since $dy_1=dy_2=0$.)

\textbf{Maslov index at the vertex.}
The Maslov index of the loop in the Lagrangian Grassmannian
$\Lambda(\bR^4, \om) \cong U(2)/O(2)$ traced by the four faces as we go
around the vertex is:
\[
  \mu(v) = \sum_{i=1}^4 \mu(F_i \to F_{i+1}),
\]
where $\mu(F_i\to F_{i+1})$ is the Maslov contribution of the crossing
at the shared edge.  For the 4-valent model above, each adjacent pair
$(F_i, F_{i+1})$ forms a positive crossing in the Maslov sense, contributing
$+1/2$ to the total.  Hence $\mu(v) = 4\cdot(1/2) = 2$.

\textbf{Surgery at the vertex.}
A local Lagrangian smoothing at $v$ is equivalent to performing a
\emph{Lagrangian surgery} (Polterovich, or Lalonde--Sikorav) at the
intersection point of two Lagrangian planes.  For the 4-valent vertex,
the four faces can be smoothed pairwise:

\medskip
\noindent\textit{Step 1}: Smooth the pair $(F_1,F_2)$ by the standard
``handle'' construction: replace $F_1\cup F_2$ near their shared edge
$\{x_1=0, x_2=0, y_2=0, y_1=0\}$ by a smooth Lagrangian annulus
$A_{12}$ that is Hamiltonian-isotopic to $F_1\cup F_2$.  This introduces
a small perturbation supported near the shared edge.

\noindent\textit{Step 2}: Similarly smooth $(F_3,F_4)$ by an annulus $A_{34}$.

\noindent\textit{Step 3}: The resulting surface near $v$ consists of two
smooth Lagrangian strips $A_{12}$ and $A_{34}$ meeting transversally at
$v$ (after a small perturbation).  A second Lagrangian surgery at their
intersection produces a single smooth Lagrangian surface near $v$.

\textbf{Justification of smoothability.}
The obstruction to performing each Lagrangian surgery lies in $\pi_1(\Lambda(2))
\cong \mathbb{Z}$ (the Maslov class).  Each surgery is possible iff the local
Maslov index at the crossing is $\pm 1$.  For the 4-valent model, each
pairwise crossing has index $+1$ (positive intersection), so both Step 1 and
Step 2 can be performed.  After Steps 1--2, the resulting crossing in Step 3
also has Maslov index $+1$ (as the sum $\mu(A_{12}) + \mu(A_{34})$ around
the remaining crossing is $2 - 1 - 1 = 0$, reduced by the two surgeries),
allowing Step 3.

The local smoothing is supported in a ball $B_\epsilon(v)$ and can be chosen
to coincide with $K$ outside $B_\epsilon(v)$.
\end{proof}

%------------------------------------------------------------------
\section{Global Smoothing}

\begin{proof}[Proof of the main theorem]
Let $K$ be a polyhedral Lagrangian surface with all vertices 4-valent.

\textbf{Step 1: Edge smoothing.}
Along each edge $e$ of $K$, exactly two faces $F_e^+, F_e^-$ meet.  The
dihedral angle between $F_e^+$ and $F_e^-$ can be any angle $\theta_e\in
(0,\pi)$.  We replace each edge with a smooth Lagrangian strip of width
$\epsilon$ interpolating between $F_e^+$ and $F_e^-$.

This is possible for any dihedral angle $\theta_e \neq 0$: the space of
Lagrangian planes containing a fixed line is connected (it is
$\mathrm{SO}(2)/\mathbb{Z}_2 \cong S^1$), so there is a smooth path of
Lagrangian planes from $T_{F_e^+}$ to $T_{F_e^-}$.  The resulting smooth
strip is Lagrangian (by construction as a graph over the edge).

\textbf{Step 2: Vertex smoothing.}
After Step 1, the surface has been smoothed along all edges.  Near each
vertex $v$, we apply Lemma~\ref{lem:local} to smooth the 4-valent intersection.
The local smoothings are supported in disjoint balls $B_\epsilon(v)$ for
$\epsilon$ small enough (since the vertex set is finite and isolated).

\textbf{Step 3: Global Lagrangian isotopy.}
The resulting surface $K_\delta$ (for small $\delta > 0$) is smooth and
Lagrangian away from the vertex neighbourhoods.  Inside each vertex ball,
it is Lagrangian by Lemma~\ref{lem:local}.  Hence $K_\delta$ is a smooth
Lagrangian surface.

The family $K_t$ for $t\in(0,1]$ is constructed by scaling $\delta = \delta(t)
\to 0$ as $t\to 0$, so that $K_t\to K_0 = K$ topologically.  Each $K_t$ is
smooth Lagrangian, and the Hamiltonian $H_t$ generating the isotopy is
supported near the edges and vertices, with $\|H_t\|_{C^0}\to 0$ as $t\to 0$.

The family $K_t$ extends continuously to $t=0$ in the topology of topological
submanifolds (the sup-norm distance $d_{\mathrm{Haus}}(K_t, K)\to 0$).
Hence $K$ admits a Lagrangian smoothing.
\end{proof}

%------------------------------------------------------------------
\section{Remarks}

\begin{remark}[Role of the 4-valent condition]
The 4-valent condition (exactly 4 faces per vertex) is essential.  For a
3-valent vertex (three faces meeting), the local Maslov count gives
$\mu(v) = 3\cdot(1/2) = 3/2$, which is not an integer, indicating a
topological obstruction to smoothing.  For 5 or more faces, additional
obstructions from the Maslov class may prevent a smooth Lagrangian filling.
The 4-valent case is the ``balanced'' case that permits smooth resolution.
\end{remark}

\begin{remark}[Relation to Legendrian knot theory]
The link $\Gamma_v$ of a 4-valent vertex is a 4-valent Legendrian graph in
$(S^3,\xi_{\mathrm{std}})$.  The smoothability theorem is related to the
existence of Legendrian surgery on 4-valent Legendrian graphs (cf.\ Etnyre--Vela-Vick).
\end{remark}

\begin{remark}[Higher dimensions]
The analogous statement for polyhedral Lagrangian hypersurfaces in $\bR^{2n}$
with $n>2$ would require understanding the local Maslov obstruction at
$(2n-2)$-dimensional vertices, which is more involved.
\end{remark}

%------------------------------------------------------------------
\section*{References}

\begin{itemize}
\item M.~Audin, F.~Lalonde, and L.~Polterovich, ``Symplectic rigidity:
  Lagrangian submanifolds,'' in \emph{Holomorphic Curves in Symplectic
  Geometry}, Birkh\"auser, 1994, pp.\ 271--321.
\item V.~Polterovich, ``The surgery of Lagrange submanifolds,''
  \emph{Geometric and Functional Analysis}, 1(2):198--210, 1991.
\item R.~Sikorav, ``Quelques propriétés des plongements lagrangiens,''
  \emph{Mémoires de la SMF}, 46:151--167, 1991.
\end{itemize}

\end{document}
